\chapter{Revisión de literatura}
\label{chap:literatura}

\noindent
La relación entre fiscalización de tránsito y siniestralidad ha sido estudiada
desde perspectivas y contextos diversos, con resultados heterogéneos. Esta
tesis se sitúa en la intersección de dos literaturas: la que relaciona el
volumen de multas con la siniestralidad, y la más amplia —y más cercana al
diseño de este trabajo— sobre disuasión (\emph{deterrence}) a través de la
presencia y capacidad de fiscalización de la autoridad.

\section{Multas y siniestralidad}

Varios trabajos documentan una asociación positiva entre historial de
infracciones y riesgo de siniestro a nivel individual. \citet{factor2014}
encuentra, con regresiones logísticas sobre datos de Israel, que quienes han
recibido multas en el pasado son más propensos a sufrir hechos de tránsito.
\citet{huajing2022} halla, para la provincia de Anhui, que conductas como el
exceso de velocidad y la conducción distraída se asocian a mayor riesgo de
siniestros fatales.

En el plano causal, \citet{luca2015} usa un diseño de variables instrumentales a
partir del programa ``Click-it-or-Ticket'' en Massachusetts y concluye que las
multas ayudan a generar conciencia entre los conductores. \citet{makowsky2011}
encuentran, para municipios de Estados Unidos con restricciones
presupuestales, que un mayor volumen de multas se asocia con menos siniestros.
\citet{mphela2013} reporta, para Botswana, que el efecto de la aplicación de la
ley sobre las muertes viales es limitado, y sugiere que elevar la edad mínima
para obtener licencia podría ser más efectivo que endurecer las sanciones
económicas.

\section{Enforcement policial y disuasión}

Un cuerpo distinto de literatura —más cercano al diseño de esta tesis— no mide
el efecto de la multa individual, sino el de la \emph{capacidad de
fiscalización} de la autoridad, sea esta medida en agentes, patrullas u
operativos. El marco teórico de referencia es \citet{becker1968}: la
probabilidad percibida de ser sancionado, y no solo la severidad de la
sanción, es lo que determina el efecto disuasivo sobre el comportamiento. Esta
distinción es relevante para esta tesis porque el diseño empírico evalúa un
cambio en la \emph{capacidad de fiscalización} (número y digitalización de
agentes autorizados para infraccionar) más que en el monto de las multas.

La evidencia empírica más cercana a este trabajo es \citet{deangelohansen2014},
quienes explotan un recorte presupuestal que redujo en 35\,\% el número de
oficiales de la patrulla carretera de Oregón en 2003. Encuentran que la caída
en infracciones emitidas estuvo acompañada de un aumento significativo en
lesiones y muertes viales, y estiman que cada \$309{,}000 dólares de gasto
adicional en patrullaje carretero previene una muerte vial. El diseño de
\citeauthor{deangelohansen2014} es, en esencia, el espejo del de esta
tesis: en vez de una reducción de agentes, aquí se explota un aumento
diferencial en la capacidad de fiscalización, y en vez de variación temporal
pura, se explota variación espacial en la exposición a la infraestructura de
fiscalización. \citet{chalfinmccrary2017}, en una revisión exhaustiva de la
literatura de disuasión criminal, documentan un patrón consistente con este
resultado: existe evidencia sólida de que el comportamiento responde a la
presencia y capacidad de la policía, pero mucho menos evidencia de que
responda a la severidad de la sanción en sí misma —una distinción que motiva
tratar las multas como \emph{mecanismo} y no como la variable de interés
central de este trabajo (Sección~\ref{sec:estrategia_outcomes}).

\section{Evidencia para la Ciudad de México}

La referencia local más cercana a esta tesis es \citet{quintero2023}, quienes
evalúan mediante un diseño de series de tiempo interrumpido controlado (CITS)
un conjunto de políticas de seguridad vial de la Ciudad de México entre 2015 y
2019. Encuentran resultados opuestos según el tipo de sanción: las políticas
de 2015 —métodos automáticos de detección de exceso de velocidad, uso del
celular o giros prohibidos, con consecuencia económica directa (fotomultas)—
se asocian a una \emph{disminución} de la mortalidad vial, mientras que la
política de 2019 —el sistema de puntos sin consecuencia económica inmediata
(Fotocívicas, Sección~\ref{sec:contexto_politica})— se asocia a un
\emph{incremento}. Este hallazgo es un antecedente directo para esta tesis:
sugiere que la naturaleza de la sanción (económica vs. administrativa) importa
tanto como su volumen, y anticipa que el vínculo entre más fiscalización y
menos víctimas no es automático en el contexto específico de la Ciudad de
México.

\section{Aporte de esta tesis}

Esta tesis se distingue de la literatura revisada en tres puntos. Primero, la
mayoría de la evidencia sobre multas y siniestralidad —incluida la versión
2024 de este mismo proyecto— se concentra en el volumen de infracciones como
variable de interés y trata a la policía como un agente pasivo que solo
multa; aquí, siguiendo más de cerca a \citet{deangelohansen2014}, el objeto de
estudio es un cambio institucional en la capacidad de fiscalización, y las
multas se usan únicamente como evidencia de que ese cambio se tradujo en más
actividad de enforcement. Segundo, la evidencia para América Latina y, en
particular, para unidades espaciales finas dentro de una misma ciudad es
escasa: \citet{quintero2023} evalúa políticas a nivel ciudad; este trabajo
explota variación \emph{dentro} de la Ciudad de México mediante exposición
espacial a la infraestructura de fiscalización. Tercero, esta tesis pone
atención explícita a los efectos de reporte en los datos administrativos de
víctimas —una amenaza a la identificación que rara vez se discute
formalmente en esta literatura— al usar fallecidos, un desenlace mucho menos
sensible a la probabilidad de que un hecho se registre formalmente, como
verificación de robustez del resultado principal
(Sección~\ref{sec:contexto_espacial} y capítulo~\ref{chap:resultados}).
